% ECONOMICS_PROBLEM: {"id": "H11-187", "title": "Building accountable state capacity", "jel_code": "H11", "source_id": "OP-0381"}
% !TeX program = xelatex
\documentclass[11pt,a4paper]{article}
\usepackage[margin=23mm]{geometry}
\usepackage{fontspec}
\setmainfont{Latin Modern Roman}
\usepackage{amsmath,amssymb,mathtools}
\usepackage{xcolor,enumitem,booktabs,longtable,array,xurl,hyperref}
\definecolor{muted}{HTML}{53636C}
\hypersetup{colorlinks=true,urlcolor=blue,linkcolor=blue}
\setlength{\parindent}{0pt}
\setlength{\parskip}{5pt}
\newcommand{\R}{\mathbb R}
\newcommand{\E}{\mathbb E}
\newcommand{\N}{\mathbb N}
\newcommand{\ind}{\mathbf 1}
\newcommand{\Var}{\operatorname{Var}}
\newcommand{\Cov}{\operatorname{Cov}}
\newcommand{\supp}{\operatorname{supp}}
\DeclareMathOperator*{\argmin}{arg\,min}
\DeclareMathOperator*{\argmax}{arg\,max}
\newcommand{\entrymeta}[3]{{\sffamily\small\color{muted}\textbf{\texttt{#1}}\quad|\quad #2\par #3\par}}
\newcommand{\Problemref}[1]{\texttt{#1}}
\newcommand{\JELref}[2]{\texttt{#2} (JEL \texttt{#1})}
\begin{document}
\section*{Building accountable state capacity}
\textbf{Problem ID:} \texttt{H11-187}\quad \textbf{JEL:} \texttt{H11}\par
% BEGIN ECONOMICS STATEMENT
\label{problem:OP-0381}
{\sffamily\small\color{muted}Original subject group: Political economy and institutions\par}
\label{definitions:problem:30007663}
\textbf{Audit classification:} Published research agenda. Full review checked. Separating politician and bureaucrat roles is explicit future work; capacity/accountability constraints and reform sequencing are editorial.
\subsection*{Definitions and precise research target}
\textbf{Complete editorial problem.} Fix polities initially possessing limited fiscal, legal, and administrative capacity. Let \(p\) be a feasible sequence of capacity and accountability reforms with a stated cost and implementation horizon. Let \(K_{ct}(p)\) be verified public-service and revenue capacity, \(A_{ct}(p)\) an accountability index defined before evaluation, and \(R_{ct}(p)\) measured coercion or elite appropriation. For horizon \(T\), estimate \(\tau_K(p)=E[K_{cT}(p)-K_{cT}(0)]\), and analogous accountability and coercion effects. The design objective is to find sequences improving capacity subject to specified lower bounds on accountability and upper bounds on repression and appropriation. These constraints encode declared social objectives. Separate the roles of elected politicians, appointed officials, and nonstate actors, because identical administrative investments may operate under different political incentives. Identification requires credible reform timing or assignment and explicit treatment of conflict, spillovers, and regime change. Evaluate service outcomes and citizens' ability to contest decisions rather than equating greater extraction with better capacity. The review explicitly requests knowledge about political oversight and bureaucratic effectiveness. Sequencing, causal estimands, and the constrained reform objective are our adaptation; the source does not offer a universally valid reform recipe.

\textbf{Definitions and admissible scope.} Capacity \(K\) is a three-component vector measuring revenue administration, legal enforcement, and service delivery, each in independently specified units. Define accountability \(A\) as a preregistered measure of contestation or oversight and repression/appropriation \(R\) as separate measured harms; higher revenue alone is not higher welfare. A policy sequence \(p=(p_0,\ldots,p_T)\) specifies budget, implementer, oversight, and support withdrawal. Use a scalar objective \(v(K)\) chosen before analysis and require \(E[A_T(p)]\ge\underline A\), \(E[R_T(p)]\le\overline R\), and \(C(p)\le B\). Average constraints differ from almost-sure safeguards and must be named. Politician and bureaucrat actions and incentives are separate primitives of any structural model; a reduced-form trial identifies the assigned bundle, not an unrestricted mechanism decomposition. Fix a baseline polity weighting law, \(T<\infty\), a feasible reform menu, and measured capacity, accountability, repression, and appropriation scales. Use \(E[v(K_T(p))]\) as the scalar capacity objective, a fixed measurable finite-mean index; \(\tau_K\) itself is a componentwise vector contrast. If repression and appropriation are separate, use two fixed harm indices and two upper bounds. Require \(E[A_T(p)]\ge\underline A\), \(E[R_T(p)]\le\overline R\), and declared expected resource cost at most \(B\). The output is the identified set of feasible reform rankings or an explicitly chosen robust optimum over the admissible causal laws. A model specifies politician, bureaucrat, and nonstate-actor actions and equilibrium selection; a bundled trial need not identify those separate actions' effects.
\subsection*{Variants and assumption changes}
\begin{enumerate}
\item Political oversight interaction (source-motivated): vary bureaucratic recruitment or incentives jointly with citizen or elected-official oversight. Estimate interactions in capacity and accountability while tracking selection into office.
\item Sequencing (editorial): compare accountability-first with capacity-first paths under equal budgets and explicit commitment rules. Require follow-up across political turnover and state whether guarantees are average or unit-specific.
\end{enumerate}
\subsection*{References and source audit}
\textbf{Support and access.} Full review checked. Separating politician and bureaucrat roles is explicit future work; capacity/accountability constraints and reform sequencing are editorial. \textbf{Locator:} Sections 4.2 and 5; conclusion p. 418, third direction.
\begin{enumerate}\raggedright
\item\hypertarget{definitions:source:30007663:1}{} Timothy Besley; Robin Burgess; Adnan Khan; Guo Xu. 2022. \emph{Bureaucracy and Development}. Section 4.2, printed pp. 413–415; concluding research agenda, printed p. 418, third direction.
\url{https://www.robinburgess.com/s/Bureaucracy_and_Development.pdf}.
DOI: \url{https://doi.org/10.1146/annurev-economics-080521-011950}.
\end{enumerate}

\subsection*{Common conventions: Source mathematical conventions}

These conventions apply to each entry unless explicitly replaced.

\textbf{Models and identification.} A primitive model \(m\) specifies all probability laws, feasible choices, information, equilibrium and measurement rules used by its target. The admissible class \(\mathcal M\) is fixed independently of outcomes. If \(P_O(m)\) is its observable probability law and \(\tau(m)\) the target, its identified set at observable law \(P\) is
\[
 \mathcal I_\tau(P)=\{\tau(m):m\in\mathcal M,\ P_O(m)=P\}.
\]
Point identification means that this set is a singleton. An empty set signals incompatibility of data and assumptions. Coordinate bounds need not describe the joint set. Bounds are sharp only if every retained target is attainable or arbitrarily approachable by compatible models. Finite bounds require explicit restrictions on unobserved quantities; unrestricted confounding, hidden wealth, extrapolation or arbitrary equilibrium selection can yield uninformative sets.

\textbf{Causal interventions.} A potential outcome \(Y_i(p)\) is the outcome under a complete policy regime \(p\), including timing, financing, information and market spillovers. The target population and units are fixed before treatment. With interference, use the complete assignment vector or a justified exposure mapping; individual no-interference notation alone cannot identify an economy-wide intervention. A difference of mean potential outcomes does not require the cross-world joint distribution. Conditioning on post-treatment migration, survival, enrollment or employment changes the estimand and needs separate justification.

\textbf{Statistical statements.} An estimator, confidence set, forecast or test specifies a sampling law, model class and loss/coverage criterion. Uniform \((1-\alpha)\)-coverage means \(\inf_{m\in\mathcal M}\Pr_m\{\tau(m)\in C_n\}\geq1-\alpha\), or an explicitly asymptotic version. The whole parameter space trivially has coverage, so informativeness requires a stated width, risk or power objective. A forecasting improvement is relative to a named benchmark, information set and evaluation population.

\textbf{Equilibrium and optimization.} An equilibrium comprises feasible plans, optimality, beliefs where needed, budget constraints and all market/resource clearing. The primitives or a stated benchmark define each correspondence \(\mathcal E(p;m)\). Nonemptiness is assumed or proved, never inferred just from compact policy sets. A selected-equilibrium target uses a declared selection rule; a robust target quantifies over all permitted equilibria. Use suprema or infima unless attainment is established. An \(\varepsilon\)-optimal policy is feasible and within \(\varepsilon\) of the finite optimum value. Report an empty feasible set separately.

\textbf{Welfare and accounting.} Normative utility, interpersonal normalization, social weights, numeraire, discounting and the use of public receipts are primitives. Behavioral utility need not coincide with the welfare criterion. Household welfare includes ownership income and rents once; adding profits again to compensating variation generally double-counts them. A monetary equivalent is defined by an explicit transfer and price convention, with monotonicity and range conditions. Average effects, mechanism contributions, transfers and welfare losses are different targets. Infinite sums require convergence and dynamic feasible plans require terminal or no-Ponzi conditions.

\textbf{Source versions.} A working-paper date, revision date and journal publication year are distinguished. A DOI for a journal version is not automatically the DOI of an earlier manuscript. Locators refer to the stated version and printed pages where specified. A metadata-only check does not verify a claim about a full-text conclusion. Every entry's source subsection narrows the provenance claim accordingly.
% END ECONOMICS STATEMENT
\end{document}
